\section{A bank and insurer share a group close}
\label{sec:sector-company-composition}

A bank's loan receipts, an insurer's claims and a parent's guarantee can all
affect the same group's cash without belonging to the same legal taxpayer or
the same shareholders. Consider a parent P that owns 80\% of bank B and
insurer I. Control and equal ordinary ownership rights are supplied facts.
P reports in USD; both subsidiaries keep their books in a foreign currency,
FCU. They apply the same retained historical EU accounting profile. Their
opening is 1 January 2025; the two intervals close on 1 January 2026 and
1 January 2027. Events described as first- or second-period events occur at
those respective closing dates. This explicit calendar also fixes which
ownership fraction receives each period's profit.

P begins with cash 500, investments 80 in each subsidiary and an external
financial-guarantee liability 6, so its equity is 654. B has cash 300,
customer deposits 200 and equity 100. I has cash 200, insurance liability 100
and equity 100. Of I's cash,20 is a demand deposit at B; the same deposit is
included in B's liability 200. The deposit is identified by its creditor I,
debtor B, currency, amount and matching closing dates. No interest or movement
occurs in it during these two periods. Opening FX is 1 USD per FCU.

The group therefore starts with cash\(500+300+200-20=980\), liabilities
\(200+100+6-20=286\), and equity 694. Eliminating the parent's investments
leaves NCI\(20+20=40\), alongside the parent's equity 654. The two component
calculators' notional funding amounts do not become new legal capital. Only
their identified loan and insurance movements enter B's and I's books.
Adding the isolated component balance sheets would answer a different question.

\subsection{Three taxpayers, three different profit origins}

At opening B advances 102 and receives an integral loan fee 2. The loan's initial
carrying amount is 100. Its receipts are 55 at the first close and 60.5 at the
second. At the supported effective interest rate 10\%,
\begin{align}
 L_1&=100(1.10)-55=55,\\
 L_2&=55(1.10)-60.5=0.
\end{align}
Interest income is 10 and 5.5. The contractual cash record separately identifies
principal receipts 50 and 52, and coupons 5 and 8.5. Thus principal receipts return
the 102 advanced, while coupons 13.5 plus the integral fee 2 equal lifetime
effective-interest income 15.5. The first coupon is 5, although that period's
effective-interest income is 10. Cash classification cannot be recovered by
renaming the interest-income line. The admitted bank policy classifies these
lending flows as operating and preserves their separate gross origins
\citep[IAS~7.14--15,31--33; IFRS~9.B 5.4.1--B 5.4.4]{eu 20231803}.

I uses the previously derived loss-component service schedule. Its opening
liability 100 incurs finance expense 10 and releases claims/risk adjustment 55,
closing at 55; the second finance expense is\(-5\), followed by releases 50,
closing at zero. Actual claims paid are 50 and 45. Insurance profit before tax
is\(-5\) and 10. Neither the loss component nor the risk-adjustment release is
another cash payment.

For P, the guarantee beneficiary is an unrelated external creditor. A supported
scope finding places this contract under IFRS~9. Its opening recognized amount
is 6. Supported service income is 2 in each period, while independently supplied
expected-credit-loss measurements are 7 and 3. Let \(U_t\) be the unamortized
initial amount and \(A_t\) that allowance. Subsequent liability is
\begin{equation}
 G_t=\max\{U_t,A_t\},\qquad U_1=4,\quad U_2=2.
 \label{eq:sector-guarantee-higher-of}
\end{equation}
This is the higher-of rule in IFRS~9.4.2.1(c), conditional on the stated
classification and measurements \citep{eu 20231803}. It gives liabilities 7
and 3. The excess above \(U_t\) is 3 and 1, respectively. Its changes,3 and
\(-2\), are impairment expense and reversal. Subtracting them from service
income 2 gives guarantee profit\(-1\) and 4. No guarantee settlement occurs
during these periods, and consolidation does not remove an external liability.

Tax now connects these movements to legal equity. The supplied law taxes book
profit at 20\%, including guarantee changes, with usable loss carryforwards.
There is separate convincing evidence supporting the losses' recovery; a
future book profit alone is not that evidence
\citep[IAS~12.34--36,46--47]{eu 20231803}. There is no loss pooling between P,
B and I. Outside-basis exceptions and absence of tax on translation are also
explicitly supported premises, not conclusions about an actual jurisdiction.

For each taxpayer let \(\ell\) denote the opening unused tax loss and \(p\)
the period's pretax profit. Under this finite supplied law, loss used is
\(u=\min\{\ell,\max(p,0)\}\), current tax is
\(0.2[\max(p,0)-u]\), and closing loss is
\(\ell-u+\max(-p,0)\). The recognized DTA is 20\% of that loss. Deferred
tax expense equals opening DTA minus closing DTA; hence the tax benefit on a
new supported loss is negative expense, without a current cash refund.

B pays tax 2 and 1.1 and earns after-tax profit 8 and 4.4. I's first loss 5 creates
DTA 1 and after-tax loss 4. Its next profit 10 uses that loss, pays tax 1 and
reverses DTA 1, leaving after-tax profit 8. P's first loss 1 creates DTA 0.2 and
after-tax loss 0.8. Its next profit 4 uses the loss, pays 0.6 and reverses 0.2,
leaving after-tax profit 3.2. Before any ownership purchase, the foreign legal
books therefore close as follows; all amounts in this table are FCU.

\begin{center}
\begin{tabular}{lrrrr}
\toprule & B, first & B, second & I, first & I, second\\\midrule
Cash &253&312.4&150&104\\
Loan &55&0&0&0\\
Deferred tax asset &0&0&1&0\\
Deposits or insurance liability &200&200&55&0\\
Equity &108&112.4&96&104\\\bottomrule
\end{tabular}
\end{center}
For example, B's first cash is\(300-102+2+55-2=253\), and
\(253+55-200=108\). I's first cash is 150, while its equity is
\(150+1-55=96\). The tax benefit belongs in equity through profit even
though it has not produced cash.

\subsection{Currency and ownership follow different dates}

The exact first-period transaction rate is 1.1 and the closing rate 1.2 USD/FCU;
the second rates are 1.25 and 1.3. Opening loan origination uses the opening
rate 1. These are stated event rates, not unexplained annual averages. Under
IAS~21.39--41, closing assets and liabilities use closing rates, income uses
transaction-date rates, and the resulting translation difference enters OCI
\citep{eu 20231803}.

With no foreign capital contributions or dividends in the period, a subsidiary
whose opening equity is \(E\), after-tax profit is \(p\), and opening,
event and closing rates are \(r_0,r_e,r_c\), has translation difference
\begin{equation}
 \mathrm{OCI}=E(r_c-r_0)+p(r_c-r_e).
 \label{eq:sector-translation-origin}
\end{equation}
To derive it, subtract translated opening equity \(Er_0\) and translated
profit \(pr_e\) from translated closing equity \((E+p)r_c\).
The first term revalues the opening net assets; the second reconciles the
two rates applied to this period's retained profit. Thus B's first translated
profit is 8.8 and OCI is\(100(0.2)+8(0.1)=20.8\). I's translated profit is
\(-4.4\), with OCI\(100(0.2)-4(0.1)=19.6\). Together with P's loss 0.8,
group profit is 3.6 and OCI 40.4.

Only after those first-period results does P buy another 10\% of B for USD 15.
Control continues. B's translated net assets are\(108(1.2)=129.6\), so the
NCI carrying amount removed is 12.96. Cash consideration exceeds that amount
by 2.04, charged directly to parent equity. No goodwill or profit is created
by this retained-control transaction \citep[IFRS~10.23,B 96]{eu 20231803}.
P's separate investment rises from 160 to 175 and cash falls from 500 to 485;
B makes no legal entry. First-period profit and OCI still use 80\% parent
ownership. Bank attribution changes to 90\% only for the second interval.

First-period NCI profit is\(0.2(8.8-4.4)=0.88\), and its OCI is
\(0.2(20.8+19.6)=8.08\). Consequently,
\begin{align}
 NCI_1&=40+0.88+8.08-12.96=36,\\
 E_{P,1}&=654+(3.6-0.88)+(40.4-8.08)-2.04=687.
\end{align}
The sum 723 also follows from\(694+3.6+40.4-15\). This second bridge tests
the equity transaction independently of the closing asset total.

In the second interval, B's translated profit is 5.5 and OCI
\(108(0.1)+4.4(0.05)=11.02\). I's translated profit is 10 and OCI
\(96(0.1)+8(0.05)=10\). Group profit is 18.7 after adding P's 3.2, and
OCI is 21.02. NCI receives\(0.1(5.5)+0.2(10)=2.55\) of profit and
\(0.1(11.02)+0.2(10)=3.102\) of OCI. Closing NCI is 41.652, with parent
equity 721.068 and total equity 762.72.

\subsection{Cash translation does not equal translation OCI}

At the first close, I's deposit 20 appears as 24 in translated cash and as 24
in B's translated deposits. The group eliminates both; neither legal book
changes. Group cash is\(485+(253+150-20)(1.2)=944.6\). Adding loan 66 and
DTAs 1.4 gives assets 1012; liabilities are
\((200+55-20)(1.2)+7=289\). At the second close the deposit elimination
is 26. Group cash becomes\(484.4+(312.4+104-20)(1.3)=999.72\), and liabilities
are\((200-20)(1.3)+3=237\). There is then no loan or DTA.

Cash must also be reconstructed from its own dated origins. First-period
operating cash is opening net advance\(-100\), plus first-period loan receipt,
bank tax and insurer claims translated at 1.1:
\begin{equation}
 CFO_1=-100+(55-2-50)(1.1)=-96.7.
\end{equation}
The NCI purchase is financing cash\(-15\), under IAS~7.42 A--42 B.
The currency effect on B's cash is
\(300(0.2)-100(0.2)+53(0.1)=45.3\); on I's cash it is
\(200(0.2)-50(0.1)=35\). The internal deposit's currency effect 4 must also
be eliminated. Thus cash FX is 76.3 and
\begin{equation}
 980-96.7-15+76.3=944.6.
\end{equation}
In the second period, operating cash is
\((60.5-1.1-45-1)(1.25)-0.6=16.15\). B's cash currency effect is 28.27,
I's 12.7, and the internal deposit's 2, giving net 38.97. The bridge is
\(944.6+16.15+38.97=999.72\). IAS~7.28 requires the exchange effect needed
to reconcile opening and closing cash separately from cash flows
\citep{eu 20231803}. It is not the OCI 40.4 or 21.02, which also reflects
noncash net assets.

Keeping all rates at 1 removes both cash FX and translation OCI, while
preserving the legal tax and operating calculations. Omitting the NCI purchase
leaves both subsidiary books unchanged, raises each subsequent group cash
balance by 15 and retains 20\% NCI participation in B's second-period results.
These comparisons reveal errors that a single balanced closing statement could
hide. The companion \code{SectorCompany/1} report preserves separate component
contexts, matched deposit identities, legal books, tax origins, dated ownership
and gross cash movements; it refuses missing support or mismatched inputs
without publishing a partial result. This finite composition supplies the
calculation just derived. Prudential capital, actual tax law, general guarantee
settlement and nonlinear insurance valuation still require their own methods.
